\documentclass[11pt,a4paper]{article}
\usepackage[utf8]{inputenc}
\usepackage[T1]{fontenc}
\usepackage{lmodern}
\usepackage[margin=0.5in]{geometry}
\usepackage{enumitem}
\usepackage{hyperref}
\usepackage{tabularx}
\usepackage{titlesec}

% Remove page numbers
\pagenumbering{gobble}

% Section formatting
\titleformat{\section}{\large\bfseries\uppercase}{}{0em}{}[\titlerule]
\titlespacing{\section}{0pt}{6pt}{4pt}

% Hyperlink setup
\hypersetup{
    colorlinks=true,
    linkcolor=black,
    urlcolor=blue
}

% Custom commands
\newcommand{\resumeSubheading}[4]{
  \item[]
    \begin{tabularx}{\textwidth}{X r}
      \textbf{#1} & #2 \\
      \textit{#3} & \textit{#4}
    \end{tabularx}
}

\setlist[itemize]{noitemsep, topsep=0pt, leftmargin=*, label=$\bullet$}

\begin{document}

% Header
\begin{center}
    {\Huge \textbf{MOHIT KUMAR}} \\
    \vspace{2pt}
    {\large Security Engineer $|$ AI for Security \& Security for AI} \\
    \vspace{3pt}
    \href{mailto:mohit.kr1103@gmail.com}{mohit.kr1103@gmail.com} $|$ 9973395247 $|$ Bengaluru, India $|$ \href{https://www.linkedin.com/in/mohitkumar111}{LinkedIn} $|$ \href{https://github.com/mk12002}{Github} $|$ \href{https://x.com/mohitkr111}{X} $|$ \href{https://mohitkumar-mu.vercel.app/}{Portfolio}
\end{center}

\vspace{-8pt}

% Education
\section{Education}
\vspace{-2pt}
\resumeSubheading
{Vellore Institute of Technology (VIT Chennai)}{Chennai, India}
{Bachelor of Technology in Electronics and Computer Engineering $|$ \textbf{CGPA: 9.12/10.0}}{Aug 2022 -- July 2026}

\resumeSubheading
{Radiant International School}{Patna, India}
{Higher Secondary Education $|$ \textbf{Class XII: 95.6\%} $|$ \textbf{Class X: 94.2\%}}{2019 -- 2021}

\vspace{-2pt}

% Experience
\section{Professional Experience}

\resumeSubheading
{ITC Infotech India Ltd.}{Bengaluru, India (On-Site)}
{Security Engineer (IS2)}{16 Jul 2026 -- Present}
\begin{itemize}
    \item Promoted to full-time Security Engineer on completion of the internship, driving \textbf{AI-for-Security and Security-for-AI} initiatives — building autonomous detection tooling and hardening ML/agent supply chains alongside day-to-day SOC and VAPT operations.
    \item Supporting \textbf{GRC} engagements including \textbf{J-SOX} and \textbf{MICS} control auditing across IT general controls (access, change, and segregation-of-duties), testing control design and effectiveness and driving remediation of identified gaps to closure.
\end{itemize}

\vspace{3pt}
\resumeSubheading
{ITC Infotech India Ltd.}{Bengaluru, India (On-Site)}
{Cybersecurity Intern}{05 Jan 2026 -- 15 Jul 2026}
\begin{itemize}
    \item Designed and deployed a multi-agent AI email security system leveraging 7 parallel ML agents, achieving ~96–97\% threat detection accuracy with <0.8\% miss rate and ~1.3\% false positives, significantly improving phishing detection reliability.
    \item Reduced SOC email triage time from ~30 minutes to a few seconds by building a real-time, explainable decision pipeline with automated threat correlation, counterfactual analysis, and SOC-ready attack narratives.
    \item Engineered a low-latency (<100 ms), high-throughput distributed architecture with ~99.9\% system uptime, enabling scalable, cost-efficient phishing detection without GPU dependency while improving analyst productivity.
\end{itemize}

\vspace{-2pt}

% Projects
\section{Projects}
\vspace{-2pt}
\href{https://github.com/mk12002/Bulwark}{\textbf{Bulwark — The Security Stack for Agentic AI}} $|$ \textit{Python, CycloneDX AI-BOM, SARIF, NIST AI RMF / EU AI Act}
\begin{itemize}
    \item Built a defensive security suite for agentic-AI supply chains as three composable, deterministic scanners over one engine: \textbf{Airlock} statically disassembles model artifacts (pickle / safetensors / GGUF / ONNX / Keras) and probes live MCP servers \textit{without ever executing them}; \textbf{Warden} imports an agent (LangChain/LangGraph, CrewAI, MCP config), builds its capability graph and scores OWASP-LLM ``excessive agency'' with a least-privilege before/after rewrite; \textbf{Manifest} emits a CycloneDX/SPDX AI-BOM folded into one governance report mapped to NIST AI RMF and the EU AI Act.
    \item On a corpus of \textbf{evasive adversarial pickle-RCE payloads, Airlock detected 14/14} — including the format/extension-confusion and Fickling-style import tricks that slip past picklescan, modelscan and fickling — with \textbf{zero false alarms} across benign models.
    \item Field-tested against \textbf{19 real public HuggingFace models}: \textbf{100\%} carried at least one supply-chain finding and \textbf{95\%} shipped executable pickle weights, quantifying how routinely production models ship code that runs on load.
\end{itemize}

\vspace{3pt}

\href{https://github.com/mk12002/Intermediate-Representation-IR-Framework-for-SIEM-Rule-Generation}{\textbf{Schema-Grounded NL$\rightarrow$KQL — An Intermediate Representation for SIEM Rule Generation}} $|$ \textit{LLMs, LangGraph, Pydantic, ASIM, Microsoft Sentinel}
\begin{itemize}
    \item Borrowing intermediate-representation design from compilers, built a pipeline that turns a natural-language detection requirement into a typed, schema-validated \textbf{Security IR} grounded in Microsoft Sentinel's ASIM schema, then \textit{deterministically compiles} it to KQL — separating ``understand the threat'' from ``remember the syntax,'' with a closed-loop repair mechanism and an ambiguity check that asks a clarifying question instead of emitting a confidently-wrong rule.
    \item Evaluated two pipelines on identical inputs (direct single-prompt generation vs.\ IR-mediated) so every difference is attributable to the IR layer: \textbf{field validity rose from 6.7\% to 86.7\%} (field hallucination cut $\sim$13$\times$, McNemar $p<$1e-7), syntax validity reached \textbf{97.8\%}, and the repair loop recovered \textbf{96.2\%} of failing cases within three attempts.
\end{itemize}

\vspace{-2pt}

% Publications
\section{Publications}
\vspace{-2pt}
\begin{itemize}
    \item \textbf{Hybrid Intensity Normalization and Translation Network for Structure-Preserving Intensity Normalization in Medical Imaging} $|$ \textit{Scientific Reports (Nature Portfolio)} $|$ \href{https://doi.org/10.1038/s41598-026-67051-6}{DOI: 10.1038/s41598-026-67051-6} \hfill 2026
    \item \textbf{HybEx-Law: A Hybrid Neural-Symbolic System for Legal Aid Eligibility Determination} $|$ \textit{Accepted --- Intl. Conf. on Computational Intelligence and Network Systems (CINS 2026), BITS Pilani Dubai} \hfill Oct 2026
    \item \textbf{AgriCure: An AWS S3-Integrated Deep Learning Platform for Automated Crop Disease Detection} $|$ \textit{IEEE Xplore (i-PACT 2025)} $|$ \href{https://ieeexplore.ieee.org/document/11307938}{Link to Paper} \hfill Dec 2025
\end{itemize}

\vspace{-2pt}

% Technical Skills
\section{Technical Skills}
\vspace{-2pt}
\textbf{Programming:} Python, Java, SQL, Prolog, Bash \\
\textbf{Security Engineering:} AI/ML \& Agent Security, Software Supply-Chain Security, Kubernetes \& RBAC Security, Post-Quantum Cryptography, DevSecOps \& CI/CD Security, SIEM \& Detection Engineering (Sentinel/KQL, Splunk), VAPT (Burp Suite, Nmap, Nessus), SOC \& Threat Hunting, GRC (J-SOX, MICS), MITRE ATT\&CK \\
\textbf{AI/ML Frameworks:} PyTorch, TensorFlow, Scikit-Learn, LangGraph, LangChain, HuggingFace, XGBoost, OpenCV \\
\textbf{Tools \& Technologies:} FastAPI, Docker, Git/GitHub, Azure AI, MongoDB, Streamlit, Jira, Jupyter \\
\textbf{Relevant Coursework:} Operating Systems, Computer Networks, DBMS, Data Structures \& Algorithms, Deep Learning, Computer Vision, NLP

% Certifications & Achievements
\section{Certifications \& Achievements}
\vspace{-2pt}
\textbf{Certifications:} \href{https://coursera.org/share/62b7320c3e53f81b4150e23d3e86c177}{IBM Cybersecurity Analyst Professional Certificate} $|$ \href{https://coursera.org/share/ad7d41c1bf4e64d3c3cffe247266a968}{Stanford Machine Learning Specialization} $|$ \href{https://learn.microsoft.com/en-us/users/mohitkumar-5451/credentials/fb7576bab685e9be}{Microsoft Azure AI Engineer Associate} $|$ \href{https://coursera.org/share/41ae000d50e2502a39c4c17e39b36ca6}{University of London Cyber Security Fundamentals} $|$ \href{https://catalog-education.oracle.com/ords/certview/sharebadge?id=CBA45575ECD2EA410B506FCD77CA2328F3CC8F54EE1EC201032A8490BDB46348}{OCI Generative AI Professional} $|$ \href{https://www.credly.com/badges/502f6d52-f56b-4bd7-8b4c-e3364c5acc8b/linked_in_profile}{Azure AI Fundamentals} \\
\textbf{Achievements:} IEEE IC Hackathon 2.0 Finalist (AIR 7/600 teams) $|$ 2nd Position State-level Science Exhibition $|$ Coordinated Bihar's first Atal Tinkering Lab (Niti Aayog funded) $|$ Languages: English, Hindi, Bengali

\end{document}
